\documentclass[11pt,a4paper]{article}

% ---------- Codifica e lingua ----------
\usepackage[T1]{fontenc}
\usepackage[utf8]{inputenc}
\usepackage[main=italian,provide=*]{babel}

% ---------- Matematica ----------
\usepackage{amsmath,amssymb}
\usepackage{mathtools}

% ---------- Impaginazione ----------
\usepackage[a4paper,margin=2.7cm]{geometry}
\usepackage{caption}
\usepackage{booktabs}
\usepackage{enumitem}
\usepackage{placeins}
\usepackage[colorlinks=true,linkcolor=black,citecolor=black,urlcolor=blue!55!black]{hyperref}

\newcommand{\ket}[1]{\lvert #1 \rangle}
\newcommand{\bra}[1]{\langle #1 \rvert}

\title{Il dimero\\
\large Panoramica della pipeline (vista a 2000 piedi)}
\author{Samuele \\ Relatore: Prof.\ Alessandro Chiesa}
\date{}

\begin{document}
\maketitle

\begin{abstract}
\noindent Questo documento è una mappa, non una derivazione: serve a
vedere in un colpo d'occhio come è strutturata l'intera pipeline sul
dimero --- sistema chiuso e poi sistema aperto --- prima di entrare nel
dettaglio di ciascuno stadio (ciascuno ha un proprio documento
dedicato). Copre entrambe le parti in un'unica narrazione continua:
Parte 1 (Sez.~\ref{sec:parte1}), il sistema chiuso, stato puro
$\ket\psi$ che evolve unitariamente; Parte 2 (Sez.~\ref{sec:parte2}),
lo stesso sistema reso aperto dal rumore, operatore densità $\rho$.
Scope, canali di rumore e decisioni operative per la Parte 2 sono
quelli confermati dal relatore (risposta del 2 settembre 2026,
\texttt{domande\_relatore.md}, sez.~12) e seguono la terminologia delle
sue slide (\texttt{7Physical\_implementation.pdf}).
\end{abstract}

\tableofcontents

\section{Parte 1 --- Il sistema chiuso}
\label{sec:parte1}

\subsection{Da dove partiamo: due spin, un legame}

Il dimero --- due spin $1/2$ accoppiati da un singolo legame di scambio
isotropo $J$, in campo $b$ --- è il sistema più semplice su cui costruire
l'intera pipeline: ogni spin si mappa direttamente su un qubit, nessuna
scelta di topologia, nessuna frustrazione geometrica (serve almeno un
ciclo dispari di tre o più siti per quello --- da cui il trimero ad
anello come passo successivo). È qui che si stabiliscono, nella forma
più semplice possibile, tutti gli ingredienti concettuali poi riusati
sui sistemi più grandi: la scelta dell'ansatz motivato dalla fisica
contro quello generico, la fattorizzazione di Trotter, il circuito
Hadamard test per i correlatori dinamici.

\subsection{Il termine DM: qui non c'è ambiguità di forma}

Con un solo legame, il termine Dzyaloshinskii--Moriya ha un'unica forma
possibile --- a differenza del trimero, dove tre legami aprono una
scelta fra combinazioni di segno inequivalenti (Opzione A/B). Il DM qui
serve a un solo scopo, ma essenziale: al campo critico $B/J=2$ (dove,
$D=0$, il fondamentale cambia bruscamente natura da singoletto a
$\ket{\!\downarrow\downarrow}$), $D\neq0$ apre l'anticrossing --- il gap
resta finito, la transizione diventa un crossover continuo invece di un
salto. È la stessa firma fisica già stabilita qui che, sul trimero,
richiede invece una scelta di forma per essere riprodotta.

\subsection{Gli stadi di Parte 1}
\label{sec:stadi-parte1}

\begin{table}[htbp]
\centering
\begin{tabular}{@{}cll@{}}
\toprule
\# & Stadio & Contenuto \\
\midrule
1 & Il sistema & spettro esatto, gap, magnetizzazione \\
2 & Il VQE & HA contro PMA (base a rami, esteso) \\
3 & La dinamica (Trotter) & due livelli, non tre \\
4 & I correlatori & 36 combinazioni \\
\bottomrule
\end{tabular}
\caption{I quattro stadi di Parte 1. Ciascuno ha un proprio documento
dedicato nella serie narrativa (Documenti 0--4).}
\end{table}
\FloatBarrier

\subsubsection{Il sistema}

Diagonalizzazione esatta della matrice $4\times4$: singoletto e i tre
stati di tripletto, energie in forma chiusa. Il risultato che conta per
il resto del lavoro: a $D=0$ il fondamentale cambia natura bruscamente a
$B/J=2$ (livelli che si incrociano esattamente); con $D\neq0$ il gap
resta aperto ovunque (minimo verificato $0.565$ a $D/J=0.2$) --- la
stessa distinzione, in forma più semplice, del blocco A/C del trimero.

\subsubsection{Il VQE}

Due famiglie di ansatz a confronto: HA (hardware-efficient, $n\_local$,
6 parametri, nessuna informazione sulla fisica del problema) contro PMA
(fisicamente motivato). Il PMA ha a sua volta due varianti: \textbf{base}
(1 parametro, CNOT--$R_y$--CNOT, stato iniziale scelto a rami --
singoletto sotto l'incrocio, $\ket{11}$ sopra) e \textbf{esteso}
(PMA-2q.3, RBS più due $R_y$ indipendenti, 3 parametri). Senza DM,
entrambi raggiungono $\mathcal F=1$. Con DM acceso: PMA base mostra un
calo di fedeltà reale (minimo $\approx0.82$) proprio all'anticrossing
--- un solo parametro e uno stato iniziale rigido non bastano a
interpolare con continuità attraverso il crossover indotto da
$D\neq0$; PMA esteso ripara il problema ($\mathcal F>0.9998$ su tutto lo
sweep), con la metà dei parametri di HA. Stessa lezione del trimero
(dove l'ansatz a rami minimo fallisce e quello esteso ripara), qui nella
sua forma più semplice, a un solo legame.

\subsubsection{La dinamica (Trotter)}

Struttura a \textbf{due} livelli, non tre come sul trimero: campo e
scambio fattorizzano esattamente in $H_1$ (nessuna approssimazione, dato
che $[S_z^\text{tot},H_\text{ex}]=0$ anche con un solo legame), il
termine DM va in $H_2$. A differenza del trimero, qui non serve un
Trotter \emph{interno} per lo scambio: un solo legame non ha bond
condivisi da spezzare ulteriormente. Punto di lavoro dedicato per la
dimostrazione ($b/J=-0.18$, $D/J=1$, stato iniziale $\ket{00}$,
deliberatamente asimmetrico --- altrimenti l'errore di Trotter non
emergerebbe): verificato che circuito vero e matrice compatta
coincidono a precisione macchina, pendenze $1/N$ (distanza fra
operatori) e $1/N^2$ (perdita di fidelity) confermate.

\subsubsection{I correlatori}

Circuito Hadamard test (ancilla + 2 qubit di registro). Con due siti e
tre componenti, le combinazioni sono $36$ (non $81$ come sul trimero).
Stesso argomento per gli zeri all'istante iniziale (12/36, stato
fondamentale reale) --- ma qui, a differenza del trimero (dove 4/81
restano strutturalmente nulli per ogni $t$), nessuna delle 36
combinazioni resta nulla per $\max_t$: gli zeri a $t=0$ sono tutti zeri
\emph{iniziali}, non strutturali. Preparazione VQE reale confrontata
con ampiezze esatte: coincidono entro $10^{-7}$--$10^{-6}$ per l'ansatz
esteso.

\section{Parte 2 --- Il sistema aperto (rumore)}
\label{sec:parte2}

\subsection{Da dove partiamo}

La parte a sistema chiuso è completa sul dimero (Sez.~\ref{sec:parte1}):
Hamiltoniana esatta,
VQE con termine DM, evoluzione temporale $e^{-iHt}$ via
Suzuki--Trotter, e correlazioni dinamiche $C_{ij}^{\alpha\beta}(t)$
tramite un circuito di tipo Hadamard test. Tutto lavora su uno stato
\emph{puro} $\ket{\psi}$ che evolve unitariamente: non c'è interazione
con l'ambiente, non c'è rumore.

Questa pipeline ripercorre esattamente la stessa catena di stadi ---
preparazione dello stato fondamentale, evoluzione temporale, misura dei
correlatori --- ma la accompagna con due sorgenti di rumore realistiche,
quelle di un dispositivo quantistico vero. Non si introduce fisica nuova
nell'Hamiltoniana: si aggiunge un modello di come un computer quantistico
reale esegue (in modo imperfetto) gli stessi circuiti già scritti.

\subsection{Cosa cambia: da \texorpdfstring{$\ket{\psi}$}{|psi>} a \texorpdfstring{$\rho$}{rho}}

Il rumore accoppia il sistema a un ambiente che non osserviamo. Uno
stato puro non basta più a descrivere il sistema una volta che
tracciamo via i gradi di libertà dell'ambiente: serve l'operatore
densità $\rho$, e l'evoluzione non è più unitaria pura ma un'operazione
quantistica più generale. Il documento sul modello di rumore
(\texttt{risultati\_modello\_\allowbreak rumore\_dimero.tex}) richiama questo
formalismo dal punto dove serve; qui basta sapere che da questo momento
lavora con $\rho$ al posto di $\ket{\psi}$, e che in Qiskit questo
corrisponde a usare
\texttt{AerSimulator(method=\allowbreak"density\_matrix")} invece del
simulatore statevector.

\subsection{Gli stadi della pipeline}
\label{sec:stadi}

\begin{table}[htbp]
\centering
\begin{tabular}{@{}cll@{}}
\toprule
\# & Stadio & Cosa cambia rispetto al sistema chiuso \\
\midrule
1 & Modello di rumore & nuovo: nessun analogo nel sistema chiuso \\
2 & VQE con termine DM, rumoroso & stesso ansatz, simulatore a $\rho$ \\
3 & Quantum simulation (Trotter), rumorosa & stesso circuito, compilato per passo \\
4 & Correlazioni dinamiche, rumorose & stesso Hadamard test, readout incluso \\
5 & Scan sui parametri di rumore & nuovo: nessun analogo nel sistema chiuso \\
\bottomrule
\end{tabular}
\caption{I cinque stadi della pipeline. Gli stadi 2--4 sono mirror dei
corrispondenti stadi a sistema chiuso; gli stadi 1 e 5 sono propri di
questa pipeline.}
\end{table}
\FloatBarrier

\subsubsection{Modello di rumore}

Il relatore ha confermato due soli canali (slide 17 del corso): errore
di gate (canale depolarizzante, separato a 1 e 2 qubit) ed errore di
lettura (readout). $T_1$/$T_2$ sono esclusi esplicitamente. Il compito
di questo stadio è tradurre dati di calibrazione reali --- pubblicati
per un chip esistente --- in un oggetto \texttt{NoiseModel} di Qiskit,
passando dal formalismo a operatori di Kraus richiamato a lezione.
Documento: \texttt{risultati\_modello\_rumore\_dimero.tex} (con
derivazione completa in \texttt{teoria\_modello\_rumore.tex}).

\subsubsection{VQE con termine DM, sotto rumore}

Si ripete la stima del ground state con lo stesso ansatz PMA-2q$\cdot$3
già validato a sistema chiuso (fedeltà 1 nel caso ideale), ma il
circuito viene ora eseguito su un simulatore a operatore densità con il
\texttt{NoiseModel} del modello di rumore agganciato. L'osservabile
d'interesse è quanto l'energia stimata e la fedeltà rispetto al ground
state esatto si allontanano dal caso ideale. Documento:
\texttt{risultati\_vqe\_dm\_rumoroso\_dimero.tex}.

\subsubsection{Quantum simulation (Trotter), sotto rumore}
\label{subsec:trotter}

Stesso circuito di Suzuki--Trotter del sistema chiuso. Qui c'è
un'insidia tecnica verificata esplicitamente: se si transpila
\emph{l'intero} circuito a un livello di ottimizzazione alto, Qiskit
riconosce che gli $N$ passi identici formano nel complesso un unico
unitario a 2 qubit e li collassa in un circuito a costo
\emph{costante}, indipendente da $N$. Questo distruggerebbe
silenziosamente proprio il compromesso che vogliamo studiare (più passi
$\Rightarrow$ meno errore di Trotter ma più rumore accumulato). La
soluzione è transpilare una volta il singolo passo al suo costo minimo
(3 CNOT, già dimostrato in
\texttt{circuito\_compatto\_teoria\_e\_limiti.tex}) e poi ripetere il
blocco già compilato $N$ volte, senza ritranspilare il circuito
completo. Documento: \texttt{risultati\_trotter\_rumoroso\_dimero.tex}
--- include anche un cross-check indipendente (circuito Qiskit contro
pura potenza di matrice in \texttt{numpy}) e il confronto diretto fra
$F(N)$ a rumore nullo e a rumore acceso, per rendere esplicito perché
esista un compromesso.

\subsubsection{Correlazioni dinamiche, sotto rumore}

Stesso circuito Hadamard test del sistema chiuso (ancilla + 2 qubit di
registro), ora con rumore su tutti i suoi ingredienti: preparazione
dello stato, blocco $U(t)$ ripetuto (stadio precedente), e infine la
misura sull'ancilla, che include anche l'errore di lettura del modello
di rumore. Documento:
\texttt{risultati\_correlatori\_rumorosi\_dimero.tex}.

\subsubsection{Scan sui parametri di rumore}

Il relatore chiede esplicitamente di non fermarsi a un solo punto: si
ripete il conto dello stadio precedente (o già di quello di Trotter) al
variare di $(\varepsilon_{1q},\varepsilon_{2q},p_\text{readout})$, per
vedere come si sposta il numero di passi ottimale $N^*$ che minimizza
l'errore totale --- risultato di un compromesso fra errore di Trotter
(diminuisce con $N$) e rumore accumulato (cresce con $N$). Documento:
\texttt{risultati\_scan\_parametri\_rumore\_dimero.tex} --- include anche
un affinamento sulla legge di scala di $N^*$ verso la saturazione.

\subsection{Le due estensioni facoltative}
\label{sec:estensioni-dimero}

Oltre ai cinque stadi principali, due estensioni non prioritarie ma
completate e verificate --- non chiudono la pipeline (che resta chiusa
sui cinque stadi, Sez.~\ref{sec:stadi}), la arricchiscono.

\paragraph{VQE noise-aware.} Invece di riusare i parametri ideali
(Stadio 2), si riottimizza l'energia \emph{direttamente sotto rumore}.
Sul dimero, a differenza del trimero (dove questa estensione trova un
vantaggio reale, $\Delta E\approx-4.16\times10^{-3}$):
\textbf{nessun vantaggio significativo} a riottimizzare --- $\Delta
E\approx-1.16\times10^{-7}$, $\Delta F\approx+1.76\times10^{-8}$,
entrambi compatibili con zero entro la precisione numerica. $N^*$ del
correlatore resta $5\to5$, invariato. Documento:
\texttt{risultati\_vqe\_noise\_aware\_definitivo.tex}.

\paragraph{Readout asimmetrico.} Il caso simmetrico ($p_{01}=p_{10}$)
usato negli Stadi 1--5 è una semplificazione: su hardware reale il
rilassamento $T_1$ durante la misura rende tipicamente $p_{10}>p_{01}$.
Verificato: $N^*=5$ invariante su un rapporto $p_{10}/p_{01}$ da
$1\times$ a $50\times$, ben oltre il range realistico --- stesso esito
qualitativo del trimero. Documento:
\texttt{risultati\_readout\_asimmetrico.tex}.

\subsection{Cosa resta fuori}

Per completezza, dichiarato esplicitamente per non lasciarlo implicito:
rilassamento e defasamento ($T_1$, $T_2$) sono esclusi per decisione
del relatore, anche se compaiono nelle sue slide come canali standard;
le reti neurali (NQS) restano fuori scope come nel sistema chiuso.

\subsection{Il filo conduttore verificato in ogni stadio}

Una proprietà usata, esplicitamente o implicitamente, in ogni stadio a
valle del modello di rumore: il readout simmetrico è sempre un fattore
\textbf{moltiplicativo} sull'osservabile misurato, mai additivo --- da
cui discende, per esempio, che $N^*$ non dipende da $p_\text{readout}$
(stadio 5). Questa proprietà smette di valere non appena il readout
diventa asimmetrico (\texttt{teoria\_readout\_\allowbreak asimmetrico.tex}): un
promemoria che ogni argomento di invarianza in questa pipeline poggia su
un'ipotesi verificata, non su una coincidenza.

\begin{quote}
\textit{Ultimo passo (addendum su stile e struttura, aggiornato):}
prodotti e verificati end-to-end anche i due pacchetti riproducibili ---
\texttt{pacchetto\_dimero\_\allowbreak riproducibile.zip} (Parte 1) e
\texttt{pacchetto\_dimero\_\allowbreak rumoroso\_\allowbreak
riproducibile.zip} (Parte 2, questa pipeline). Il pacchetto di Parte 2
copriva già, fin dalla prima versione, entrambe le estensioni
(Sez.~\ref{sec:estensioni-dimero}) --- aggiornato con 6 figure in più
(due diagrammi di circuito, il confronto su griglia estesa, lo scan
completo delle 36 combinazioni sotto rumore), nessun problema aperto.
Il pacchetto di Parte 1 aggiornato con 15 figure in più (il gap
esplicito, la riparazione a 3 parametri, il confronto circuito/matrice
per Trotter, lo scan delle 36 combinazioni, tutti i nove diagrammi di
circuito, il correlatore con preparazione VQE reale). Una discrepanza
trovata durante il lavoro ($\mathcal F=0.922857$ nel Documento 2 contro
$0.851664$ ricostruito) risolta: non un errore, ma una convenzione di
fedeltà diversa ($\lvert\langle\cdot\rvert\cdot\rangle\rvert$ invece di
$\lvert\langle\cdot\rvert\cdot\rangle\rvert^2$) isolata a un solo modulo
(\texttt{vqe\_dimer.py}) --- uniformata al quadrato (la convenzione di
Qiskit, di \texttt{vqe\_test2.py} e di tutto il trimero), con
\texttt{dimero\_02\_vqe.tex} e \texttt{dimero\_03\_dinamica.tex}
aggiornati di conseguenza (5 punti). Nessun problema aperto residuo in
nessuno dei due pacchetti.
\end{quote}

\end{document}
